\documentclass[8.2pt,a4paper,twocolumn]{extarticle}
\usepackage[utf8]{inputenc}
\usepackage[T1]{fontenc}
\usepackage[romanian]{babel}
\usepackage[margin=0.85cm,columnsep=0.5cm,top=0.7cm]{geometry}
\usepackage{amsmath,amssymb,mathtools}
\usepackage{enumitem}
\usepackage{titlesec}
\usepackage{parskip}
\usepackage[most]{tcolorbox}
\usepackage{xcolor}

\newtcolorbox{res}{colback=red!6,colframe=red!65!black,boxrule=0.6pt,arc=1.5pt,
  left=3pt,right=3pt,top=1pt,bottom=1pt,breakable}

\setlist{nosep, topsep=0.5pt, itemsep=0.5pt, leftmargin=1.0em}
\titlespacing*{\section}{0pt}{3pt}{0.5pt}
\titlespacing*{\subsection}{0pt}{1.5pt}{0pt}
\renewcommand{\baselinestretch}{0.90}
\setlength{\parindent}{0pt}
\setlength{\parskip}{1.5pt}
\pagestyle{empty}

\begin{document}
\twocolumn[
  \begin{center}
  {\large\bfseries Formule rezolvare SAE --- 2026}
  \end{center}
  \vspace{0.25cm}
]

\section*{1. Discretizare $G(s)\to G_d(z)$}
$G_{\text{metodă}}(z)=G(s)\big|_{s=f(z)}$:
\begin{res}
$$\color{blue}\text{MDS (înapoi)}:\ s\to\frac{z-1}{zT_e}\qquad
\text{MDF (înainte)}:\ s\to\frac{z-1}{T_e}$$
$$\color{blue}\text{MT (Tustin)}:\ s\to\frac2{T_e}\cdot\frac{z-1}{z+1}$$
\end{res}

\textbf{Metoda matched} (conservă poli/zerouri):
\begin{res}
Pas 1: $p_d=e^{-p_cT_e}$, $z_d=e^{-z_cT_e}$ (poli \& zerouri finite)

Pas 2 (opțional): dacă $n>m$ ($n$=grad numitor, $m$=grad numărător): se adaugă $\boldsymbol{(n-1)}$ \textbf{factori} $(z+1)$ la numărător (independent de $m$!)

Pas 3: $\displaystyle\lim_{s\to0}s^\sigma G(s)=\lim_{z\to1}\Big(\frac{z-1}{zT_e}\Big)^{\!\sigma}\!G_d(z)$, $\sigma$=tipul lui $G(s)$ (nr. poli în $s=0$)
\end{res}
Ex. $G(s)=\dfrac1{s(s+b)}$ ($n=2,m=0,\sigma=1$): se adaugă 1 factor $(z+1)$:
$$G_d(z)=\frac{K_d(z+1)}{(z-1)(z-e^{-bT_e})}$$

\textbf{Circuit închis, reacție unitară:} $H(z)=\dfrac{G(z)}{1+G(z)}$.

Stabilitate: poli $H(z)$ în interiorul cercului unitar $|z|<1$.

\section*{2. Eroare staționară}
$$e_{ss}=\lim_{k\to\infty}e[k]=\lim_{z\to1}(1-z^{-1})E(z),\quad E(z)=\frac{R(z)}{1+G_d(z)}$$
\begin{res}
$$K_p=\lim_{z\to1}G_d(z)\qquad e_{ss}(\text{treaptă})=\frac1{1+K_p}$$
$$K_v=\frac1{T_e}\lim_{z\to1}(z-1)G_d(z)\qquad e_{ss}(\text{rampă})=\frac1{K_v}$$
$$K_a=\frac1{T_e^2}\lim_{z\to1}(z-1)^2G_d(z)\qquad e_{ss}(\text{parabolă})=\frac1{K_a}$$
\end{res}
\textbf{Tip sistem} = nr. poli ai $G_d(z)$ în $z=1$:
\begin{itemize}
\item tip 0 (niciun pol în $z=1$): $K_p$ \textbf{finit} $\Rightarrow e_{ss}$(treaptă) finită $\ne0$; $K_v=0\Rightarrow e_{ss}$(rampă)$=\infty$.
\item tip 1 (1 pol în $z=1$): $K_p=\infty\Rightarrow e_{ss}$(treaptă)$=0$; $K_v$ finit $\Rightarrow e_{ss}$(rampă) finită $\ne0$; $K_a=0\Rightarrow e_{ss}$(parabolă)$=\infty$.
\end{itemize}

\section*{3. Transformata Z inversă}
$Y(z)=G_{er}(z)R(z)$, $F(z)=B(z)/A(z)$, poli simpli $p_i$:
\begin{res}
Fracții simple (dacă $Y(z)$ improprie, se împarte prin $z$ întâi):
$$\frac{Y(z)}{z}=\sum_i\frac{c_i}{z-p_i},\quad c_i=(z-p_i)\frac{Y(z)}{z}\bigg|_{z=p_i}$$
$$\Rightarrow y[k]=\sum_i c_ip_i^k$$
\end{res}
Poli complecși $p_{1,2}=\alpha\pm j\beta=re^{\pm j\theta}$ ($r=\sqrt{\alpha^2+\beta^2}$, $\theta$ cu semn de cadran): se grupează termenii conjugați folosind
$$\mathcal Z\{r^k\cos k\theta\}=\frac{z(z-r\cos\theta)}{z^2-2rz\cos\theta+r^2}$$
$$\mathcal Z\{r^k\sin k\theta\}=\frac{zr\sin\theta}{z^2-2rz\cos\theta+r^2}$$

Tabel transformate uzuale ($k\ge0$):
\begin{center}
\small
\begin{tabular}{lll}
semnal & $F(s)$ & $F(z)$\\\hline
$\delta[k]$ & $1$ & $1$\\
treaptă & $1/s$ & $\dfrac{z}{z-1}$\\
rampă $t$ & $1/s^2$ & $\dfrac{T_ez}{(z-1)^2}$\\
$e^{-at}$ & $\dfrac1{s+a}$ & $\dfrac{z}{z-e^{-aT_e}}$\\
\end{tabular}
\end{center}
Impuls Dirac: $R(z)=1\Rightarrow Y(z)=G_{er}(z)$ direct.

\textbf{Verificare (câștig static):} $\displaystyle\sum_k y[k]=Y(z)\big|_{z=1}=H(1)$; pentru intrare impuls, $H(1)=1-e_{ss}(\text{treaptă})$ când sistemul are câștig static unitar.

\section*{4. Caracteristici de frecvență}
\subsection*{Transformarea $w$ (Bode)}
\begin{res}
$$\color{blue}z=\frac{1+\frac{T_e}2w}{1-\frac{T_e}2w}=\frac ND,\ N=1+\frac{T_e}2w,\ D=1-\frac{T_e}2w$$
\end{res}
Se înlocuiește $z=N/D$ în $G_d(z)$, se amplifică cu $D^{\deg}$ pt. a scăpa de fracții $\Rightarrow G_d(w)$. Se pune $w=j\omega$ și se lucrează direct cu $\omega$ (rad/s) ca variabilă de frecvență a graficului --- fără conversie înapoi la $\omega$.

\textbf{Poli reali + zerouri (asimptote Bode clasice):} se factorizează $G_d(w)=K_B\prod(1\pm w/w_i)$; fiecare pol/zero $w_i$ = frângere la $\omega=|w_i|$ ($-20$dB/dec per pol, $+20$dB/dec per zero); nivel start $20\lg K_B$; zerou cu $w_z>0$ (în semiplanul drept) = \textbf{neminim-fazic}: contribuie $+20$dB/dec la modul dar $-90^\circ$ la fază (ca un pol!).

\textbf{Poli complecși (formă canonică ordin 2):}
$$\color{blue}H(w)=\frac{K}{1+2\alpha Tw+T^2w^2}\ (\text{identifici }K,T,\alpha\text{ termen cu termen})$$
$\omega_f=1/T$ (frângere, $K_{dB}=20\lg K$); pantă $-40$dB/dec după $\omega_f$; fază: $0^\circ\to-90^\circ$ la $\omega=\omega_f\to-180^\circ$.

\subsection*{Curbă polară}
Se calculează \textbf{la fel ca la Bode}: substitui $z=N/D$ (transformarea $w$ de mai sus) în $G_d(z)$ $\Rightarrow G_d(w)$; apoi, în loc de a rămâne în $w$, pui direct $w=j\omega$:
\begin{res}
$$G_d(w)\big|_{w=j\omega}=G_d(j\omega)=\text{Re}(\omega)+j\,\text{Im}(\omega),\quad \omega\in[0,\infty)$$
Plecare ($\omega=0$, i.e. $z=1$): $G_d(w{=}0)$ --- real. Sosire ($\omega\to\infty$, i.e. $z=-1$): $\displaystyle\lim_{\omega\to\infty}G_d(j\omega)$ --- real (raportul coeficienților dominanți).
\end{res}
Pași: 1) $G_d(w)$ deja calculat la Bode (numărător/numitor polinoame în $w$); 2) pui $w=j\omega$; 3) amplifici cu conjugata numitorului pt. Re/Im; 4) evaluezi la capete + puncte intermediare; 5) trasezi.

\textit{Nu se substituie $z=e^{j\omega T_e}$ direct --- se trece obligatoriu prin $G_d(w)$, exact ca la caracteristica amplitudine-fază.}

\section*{5. Locul rădăcinilor în $z$}
$GH(z)=K\dfrac{(z-z_1)\cdots}{(z-p_1)\cdots}$, ec. caracteristică $1+GH(z)=0$:
\begin{res}
condiția modul: $|GH(z)|=1$\qquad condiția argument: $\angle GH(z)=180^\circ(2k+1)$
\end{res}
Nr. ramuri $=\max(n,m)$; pleacă din poli ($k=0$), se termină în zerouri finite/$\infty$ ($k\to\infty$); axa reală: la stânga unui \textbf{nr. impar} de poli+zerouri reale.
\begin{res}
Asimptote: $n-m$; $\theta_k=\dfrac{180^\circ(2k+1)}{n-m}$, $\sigma_a=\dfrac{\sum p_i-\sum z_j}{n-m}$

Ramificație: $k(z)=-\dfrac{N(z)}{D(z)}$, rădăcini reale ale $\dfrac{dk}{dz}=0\Leftrightarrow N'D-ND'=0$
\end{res}
\begin{res}
\textbf{2 poli reali + 1 zero real ($GH(z)=K\dfrac{z-z_0}{(z-p_1)(z-p_2)}$):} porțiunea complexă e mereu un \textbf{cerc centrat în $z_0$}:
$$z=z_0+w,\quad |w|^2=(z_0-p_1)(z_0-p_2)=R^2$$
Breakaway/break-in $=z_0\pm R$: coincid cu rădăcinile lui $dk/dz=0$ de mai sus (\textbf{verificare încrucișată}, nu înlocuiește calculul lor). Dacă $R>1$: cercul iese din cercul unitar $\Rightarrow$ există $k_{max}$ peste care sistemul devine instabil.

\textbf{Verificare rapidă:} $w_{b1}=z_{b1}-z_0$, $w_{b2}=z_{b2}-z_0$ (cu $z_{b1},z_{b2}$ de la ramificație) trebuie să dea $|w_{b1}|=|w_{b2}|=R$.
\end{res}

\section*{6. Model cu variabile de stare}
$G(z)=\dfrac{b_{n-1}z^{n-1}+\cdots+b_0}{z^n+a_{n-1}z^{n-1}+\cdots+a_0}$ (strict propriu):
\begin{res}
Formă complet \textbf{controlabilă}:
$$A=\begin{pmatrix}0&1&\cdots&0\\\vdots&&\ddots&\vdots\\0&0&\cdots&1\\-a_0&-a_1&\cdots&-a_{n-1}\end{pmatrix}\!,\ B=\begin{pmatrix}0\\\vdots\\0\\1\end{pmatrix}$$
$$C=\begin{pmatrix}b_0&\cdots&b_{n-1}\end{pmatrix}$$
\end{res}
Controlabilitate $P=[B\ AB\ \cdots]$: $\det P=\pm1$ \textbf{mereu} $\Rightarrow$ controlabilă prin construcție, orice $a_i$.

Observabilitate $Q=[C;CA;\ldots]^T$: observabil $\Leftrightarrow\det Q\ne0$ --- \textbf{trebuie calculat}, nu e automat.

Formă complet \textbf{observabilă} (dual prin transpunere): $A_o=A^T$, $B_o=C^T$, $C_o=B^T$. $\det Q_o=\pm1$ \textbf{mereu} (observabilă prin construcție); controlabilitatea ei ($P_o$) trebuie calculată.

Caz \textbf{biproper} ($\deg\text{num}=\deg\text{numit}=n$): $B(z)=b_nA(z)+B'(z)\Rightarrow D=b_n$, $b_i'=b_i-b_na_i$ (înlocuiesc $b_i$ în $C$ resp. $B_o$ de mai sus).

\section*{Reguli generale de reținut}
\begin{itemize}[leftmargin=1.1em]
\item stabilitate discretă: pol în \textbf{interiorul} cercului unitar $|z|<1$ (nu semiplan stâng ca la $s$); $z=e^{sT_e}$ trimite axa $j\omega\to$ cercul $|z|=1$
\item modul/unghi factor $a+jb$: $\sqrt{a^2+b^2}$, $\arctan(b/a)$ \textbf{cu semn de cadran} (dacă $a<0,b>0$: $\theta=180^\circ-\arctan(b/|a|)$)
\item la derivata $dk/dz=0$: dacă numitorul e $z^2-a_1z+a_0$ (cu minus), derivata e $2z-a_1$ (păstrează semnul) --- o inversare aici schimbă discriminantul și poate transforma rădăcini reale în complexe
\item verificare rapidă răspuns la impuls: $\sum_k y[k]=H(1)$ trebuie să coincidă cu $1-e_{ss}(\text{treaptă})$ pentru câștig static unitar
\item rotunjești numerele intermediare independent unul de altul $\Rightarrow$ mici inconsistențe la 2-3 zecimale sunt normale, nu semn de eroare
\end{itemize}

\section*{Instrucțiuni (pentru asistentul AI)}
\begin{res}
Ți se dă o poză cu un bilet de examen SAE și acest document de formule. Rezolvă biletul folosind \textbf{exclusiv} formulele de mai sus (nu introduce alte metode/notații).

\textbf{Format obligatoriu al răspunsului} --- identic cu o rezolvare oficială:
\begin{itemize}[leftmargin=1.1em]
\item Doar formule și calcule cu numerele din enunț, direct, fără nicio propoziție explicativă, fără introduceri, fără concluzii, fără reformularea cerinței.
\item Nu sări niciun pas de calcul --- fiecare linie de calcul trebuie să apară, până la rezultatul final al fiecărui subpunct.
\item Fiecare rezultat final al unui subpunct scris clar separat (ex. printr-o linie proprie sau chenar).
\item Rotunjește toate numerele afișate la maxim 2 zecimale (excepție: dacă rotunjirea la 2 zecimale ar face un număr mic, nesemnificativ, să pară 0.00 și ar induce în eroare --- atunci păstrează 3--4 zecimale doar acolo).
\item Tratează subpunctele/problemele în ordinea din enunț, fiecare separat, fără reluări sau recapitulări.
\item Adaugă grafice (folosind un instrument de plotare disponibil) oriunde subiectul cere trasare: răspuns la impuls/treaptă (stem plot), caracteristici Bode (log $\omega$, cu asimptote), curbă polară (Re vs. Im), locul rădăcinilor (plan complex, cu cercul unitate punctat).
\item Nu adăuga teorie, definiții sau explicații suplimentare --- doar rezolvarea propriu-zisă.
\end{itemize}
\end{res}

\end{document}
